\documentclass[12pt]{article}
\usepackage[utf8]{inputenc}
\usepackage[T1]{fontenc}
\usepackage[margin=1in]{geometry}
\usepackage{setspace}
\usepackage{xcolor}
\usepackage[colorlinks=true, urlcolor=blue!60!black, linkcolor=blue!60!black]{hyperref}
\setlength{\parindent}{2em}
\setlength{\parskip}{1.25em}
\renewcommand{\baselinestretch}{1.0}
\begin{document}
\begin{center}
{\LARGE \textit{Milliken v.\ Bradley}, 418 U.S.\ 717 (1974)\par}
\vspace{0.8em}
{\large Case Brief --- Four-Page Edition\par}
\end{center}
\vspace{1em}

\section{Memory Jogger}

Federal courts may not impose an inter-district metropolitan school desegregation remedy across suburban district lines unless plaintiffs prove an inter-district violation that produced a significant segregative effect in another district (\href{https://supreme.justia.com/cases/federal/us/418/717/}{Justia}).

\section{Detailed Case Facts}

In 1970 the Detroit Public Schools were about 63.6\% Black amid an overwhelmingly white tri-county metropolitan area. When the Detroit Board adopted a partial high school desegregation plan (the April 7, 1970 Plan), the Michigan Legislature passed Act 48, which effectively nullified it and mandated ``free choice'' and ``neighborhood schools.'' The NAACP and Detroit parents and students, including Ronald Bradley, sued Governor Milliken, state officials, and the Detroit Board. After a 41-day trial, the district court found extensive \textit{de jure} segregation by the Detroit Board (optional zones letting white students escape Black schools, boundaries drawn to maximize separation, one-way busing of Black students, and segregative school siting) and held the State responsible through Act 48, construction oversight, and discriminatory denial of transportation aid. It took no proof that any of the 85 suburban districts had committed segregative acts or that district lines were drawn by race.

\section{Procedural History}

Judge Stephen J.\ Roth rejected three Detroit-only plans as leaving schools 75\% to 90\% Black and accelerating white flight, declared that ``[s]chool district lines are simply matters of political convenience,'' and ordered a desegregation area of Detroit plus 53 suburban districts (about 780,000 students), with state-funded buses. The en banc Sixth Circuit affirmed liability and the need for a metropolitan remedy, but remanded so suburban districts could be joined. The Supreme Court reversed on July 25, 1974.

\section{Judicial Votes}

Reversed and remanded for a Detroit-only decree, \textbf{5--4}.
\begin{itemize}
\item \textbf{Majority:} Chief Justice Burger, joined by Justices Stewart, Blackmun, Powell, and Rehnquist.
\item \textbf{Concurrence:} Justice Stewart.
\item \textbf{Dissents:} Justice Douglas; Justice White, joined by Douglas, Brennan, and Marshall; Justice Marshall, joined by Douglas, Brennan, and White.
\end{itemize}

\section{Holding}

Under \textit{Swann}, ``the scope of the remedy is determined by the nature and extent of the constitutional violation.'' A court may not impose a multi-district remedy for single-district segregation unless discriminatory acts of the State or of local districts were a substantial cause of inter-district segregation, or district lines were deliberately drawn by race. Desegregation does not require any particular racial balance. Burger's reflections on local autonomy and on the administrative problems of consolidation are supporting rationale; the Court did not reach the suburbs' due process claims or housing-segregation claims.

\section{Analysis of Opinions}

\textbf{Majority (Burger).} The lower courts treated metropolitan racial balance as a constitutional mandate, rejecting Detroit-only plans merely because Detroit was predominantly Black. District lines are not mere ``political convenience'': local boards control taxes, bonds, curricula, and hiring, and consolidating 54 districts would make the judge a ``legislative authority'' and ``school superintendent.'' Even assuming Michigan was responsible for Detroit's violations, the record showed no cross-district segregative effect, and one 1950s tuition arrangement could not justify sweeping in 52 uninvolved districts.

\textbf{Stewart (concurring).} The holding is narrow. The city-suburb divide stemmed from ``unknown and perhaps unknowable factors,'' but inter-district relief would be proper if officials redrew lines by race or engaged in ``purposeful, racially discriminatory use of state housing or zoning laws.''

\textbf{Douglas (dissenting).} Michigan's schools are a state enterprise; combined with \textit{Rodriguez}, the ruling leaves Black students in schools ``not only `separate' but `inferior'.''

\textbf{White (dissenting).} The State, bound by the Fourteenth Amendment, cannot hide behind its own subdivisions. A metropolitan plan was ``physically easier and more practicable,'' requiring fewer buses than a Detroit-only plan.

\textbf{Marshall (dissenting).} ``[T]he Court today takes a giant step backwards.'' Roth sought only to cure proven segregation; Michigan's system was centralized; Detroit's segregation rippled across district lines; and the decision reflected ``a perceived public mood that we have gone far enough.''

\section{Examples: Future Applications}

\begin{enumerate}
\item \textbf{Same-Side: Intra-City Gerrymandering.} A city board gerrymandered zones, but neighboring suburbs have neutral borders. Relief, including compensatory programs approved in \textit{Milliken II} (1977), stays inside the city.
\item \textbf{Same-Side: ``Desegregative Attractiveness.''} Salary hikes and magnet upgrades designed to lure suburban students exceed the remedial ceiling, as \textit{Missouri v.\ Jenkins} (1995) reinforced.
\item \textbf{Opposite-Side: Secession to Preserve White Schools.} A white municipality carves out a new district to impede county-wide desegregation; under \textit{Emporia} and \textit{Scotland Neck}, cross-district relief is available.
\item \textbf{Opposite-Side: State and Suburban Discrimination.} State officials and suburbs exclude Black transfer students and site minority housing in the city, establishing an inter-district violation.
\item \textbf{Fence-Sitter: Housing and Zoning History.} Suburban redlining and racial zoning by state housing authorities, with facially unitary schools, may meet Stewart's housing exception if causation is proved.
\end{enumerate}

\section{Critique}

\textbf{Restraint.} Defenders say equity cures only proven wrongs, so binding 53 districts never found liable offends federalism and democratic accountability; local control safeguards voting and limits Article~III power. \textbf{Progressive and historical.} Critics see the turning point that halted northern desegregation: ``local control'' ignored Michigan's centralized structure; Myron and Gary Orfield stress the artificial wall between school and housing segregation built by FHA redlining and covenants; and paired with \textit{Rodriguez}, the decision locked urban students into underfunded, isolated districts and gave white flight a safe harbor.

\section{Key Quotations}

\begin{quote}
``No single tradition in public education is more deeply rooted than local control over the operation of schools; local autonomy has long been thought essential both to the maintenance of community concern and support for public schools and to quality of the educational process.''
\end{quote}
\noindent --- Chief Justice Burger, for the Court, 418 U.S.\ at 741--42.

\begin{quote}
``Before the boundaries of separate and autonomous school districts may be set aside by consolidating the separate units for remedial purposes or by imposing a cross-district remedy, it must first be shown that there has been a constitutional violation within one district that produces a significant segregative effect in another district. Specifically, it must be shown that racially discriminatory acts of the state or local school districts, or of a single school district have been a substantial cause of inter-district segregation.''
\end{quote}
\noindent --- Chief Justice Burger, for the Court, 418 U.S.\ at 744--45.

\begin{quote}
``We deal here with the right of all of our children, whatever their race, to an equal start in life and to an equal opportunity to reach their full potential as citizens. Those children who have been denied that right in the past deserve better than to see fences thrown up to deny them that right in the future. Our Nation, I fear, will be ill-served by the Court's refusal to remedy separate and unequal education, for unless our children begin to learn together, there is little hope that our people will ever learn to live together.''
\end{quote}
\noindent --- Justice Marshall, dissenting, 418 U.S.\ at 783.

\end{document}
